% GENERATED FILE. Edit canonical lessons, then run npm run generate:books.
% canonical-source-hash: fnv1a64:bedddbf846b75f21
% canonical-lessons: ZH-C376-huodong, ZH-C376-shegong, ZH-C376-aihu, ZH-C376-rangzuo, ZH-C376-guangchangwu

\chapter{Out in Public}
\label{ch:zh-out-in-public}
\hlchaptermodality{376}

\begin{chapteropening}
\textbf{By the end of this chapter:} \emph{I can talk about events, social workers, looking after the park, giving up a seat and square dancing.}

Complete the last lesson of chapter 376: I can talk about events, social workers, looking after the park, giving up a seat and square dancing.
\end{chapteropening}

\section[huódòng]{huódòng — an activity; an organised event}

\label{lesson:ZH-C376-huodong}



\begin{quote}
Before the new one: say the Mandarin for a volunteer, then the Mandarin for each other.
\end{quote}

\subsection*{You'll want to know: huódòng}
\emph{huódòng} — \textquotedblleft{}an activity; an organised event\textquotedblright{}.

\begin{quote}
\emph{zhège zhōumò gōngyuán lǐ yǒu shénme huódòng} — What's on in the park this weekend?
\end{quote}

Its characters come later; for now it is a word for the ear.

\begin{cousinweb}
\emph{huódòng} is built from \emph{huó} 'alive' and \emph{dòng} 'to move', both inherited from Middle Chinese, \emph{hwat} and \emph{duwngX} (Guangyun, in Baxter's transcription), which glosses the first 'not dead' and the second 'to stir'. The sense 'an event' is later and not given there.
\end{cousinweb}

\begin{grammarlens}[title={What you've built}]
One more word for giving reasons.
\end{grammarlens}

\subsection*{Your turn}
\begin{itemize}
\raggedright
  \item \emph{Say it:} \emph{huódòng}
  \item \emph{Say it:} \emph{huódòng}, once more
  \item \emph{Say it:} \emph{zhège zhōumò gōngyuán lǐ yǒu shénme huódòng} — What's on in the park this weekend?
  \item \emph{Recall:} say the Mandarin for a volunteer, then the Mandarin for each other, then say \emph{huódòng} again
\end{itemize}

\begin{tcolorbox}[breakable,colback=teal!4,colframe=teal!35!black,title={Before you move on}]
What does huódòng mean? (\textquotedblleft{}An activity; an organised event\textquotedblright{}.) Say it once more.
\end{tcolorbox}
\section[\texorpdfstring{\zh{社工} (shègōng)}{社工 (shègōng)}]{\zh{社工} (shègōng) — social worker}

\label{lesson:ZH-C376-shegong}



\begin{quote}
Before the new one: say the Mandarin for each other, then the Mandarin for an organised event.
\end{quote}

\subsection*{You'll want to know: \zh{社工}}
\textbf{\zh{社工}} — \emph{shègōng} — \textquotedblleft{}social worker\textquotedblright{}.

\begin{quote}
\textbf{\zh{我觉得社工的工作很重要}} \emph{wǒ juéde shègōng de gōngzuò hěn zhòngyào} — I think social workers do important work.
\end{quote}

\begin{cousinweb}
\textbf{\zh{社工}} (\emph{shègōng}) is built from \zh{社} 'society' and \zh{工} 'work', both inherited from Middle Chinese, \emph{dzyaeX} and \emph{kuwng} (Guangyun, in Baxter's transcription), which lists the first with the god of grain and glosses the second 'a craft'. The word is a modern short form of 'social work'.
\end{cousinweb}

\begin{grammarlens}[title={What you've built}]
One more word for giving reasons.
\end{grammarlens}

\subsection*{Your turn}
\begin{itemize}
\raggedright
  \item \emph{Say it:} \emph{shègōng}
  \item \emph{Say it:} \emph{shègōng}, once more
  \item \emph{Say it:} \emph{wǒ juéde shègōng de gōngzuò hěn zhòngyào} — I think social workers do important work.
  \item \emph{Recall:} say the Mandarin for each other, then the Mandarin for an organised event, then say \emph{shègōng} again
\end{itemize}

\begin{tcolorbox}[breakable,colback=teal!4,colframe=teal!35!black,title={Before you move on}]
What does \zh{社工} mean? (\textquotedblleft{}Social worker\textquotedblright{}.) Say it once more.
\end{tcolorbox}
\section[àihù]{àihù — to cherish and protect; to treat with care}

\label{lesson:ZH-C376-aihu}



\begin{quote}
Before the new one: say the Mandarin for an organised event, then the Mandarin for a social worker.
\end{quote}

\subsection*{You'll want to know: àihù}
\emph{àihù} — \textquotedblleft{}to cherish and protect; to treat with care\textquotedblright{}.

\begin{quote}
\emph{wǒmen yīnggāi àihù gōngyuán lǐ de huācǎo} — We should look after the plants in the park.
\end{quote}

Its characters come later; for now it is a word for the ear.

\begin{cousinweb}
\emph{àihù} is built from \emph{ài} 'to cherish' and \emph{hù} 'to protect', both inherited from Middle Chinese, \emph{'ojH} and \emph{huH} (Guangyun, in Baxter's transcription), which glosses the first 'to pity' and the second 'to save' and 'to help'.
\end{cousinweb}

\begin{grammarlens}[title={What you've built}]
One more word for giving reasons.
\end{grammarlens}

\subsection*{Your turn}
\begin{itemize}
\raggedright
  \item \emph{Say it:} \emph{àihù}
  \item \emph{Say it:} \emph{àihù}, once more
  \item \emph{Say it:} \emph{wǒmen yīnggāi àihù gōngyuán lǐ de huācǎo} — We should look after the plants in the park.
  \item \emph{Recall:} say the Mandarin for an organised event, then the Mandarin for a social worker, then say \emph{àihù} again
\end{itemize}

\begin{tcolorbox}[breakable,colback=teal!4,colframe=teal!35!black,title={Before you move on}]
What does àihù mean? (\textquotedblleft{}To cherish and protect; to treat with care\textquotedblright{}.) Say it once more.
\end{tcolorbox}
\section[ràngzuò]{ràngzuò — to give up one's seat (to someone)}

\label{lesson:ZH-C376-rangzuo}



\begin{quote}
Before the new one: say the Mandarin for a social worker, then the Mandarin for to treat with care.
\end{quote}

\subsection*{You'll want to know: ràngzuò}
\emph{ràngzuò} — \textquotedblleft{}to give up one's seat (to someone)\textquotedblright{}.

\begin{quote}
\emph{zài dìtiě lǐ niánqīngrén yīnggāi ràngzuò} — On the underground, young people should give up their seats.
\end{quote}

Its characters come later; for now it is a word for the ear.

\begin{cousinweb}
\emph{ràngzuò} is built from \emph{ràng} 'to yield' and \emph{zuò} 'a seat', both inherited from Middle Chinese, \emph{nyangH} and \emph{dzwaH} (Guangyun, in Baxter's transcription), which glosses the first 'to give way' and the second 'a bed or a seat'.
\end{cousinweb}

\begin{grammarlens}[title={What you've built}]
One more word for giving reasons.
\end{grammarlens}

\subsection*{Your turn}
\begin{itemize}
\raggedright
  \item \emph{Say it:} \emph{ràngzuò}
  \item \emph{Say it:} \emph{ràngzuò}, once more
  \item \emph{Say it:} \emph{zài dìtiě lǐ niánqīngrén yīnggāi ràngzuò} — On the underground, young people should give up their seats.
  \item \emph{Recall:} say the Mandarin for a social worker, then the Mandarin for to treat with care, then say \emph{ràngzuò} again
\end{itemize}

\begin{tcolorbox}[breakable,colback=teal!4,colframe=teal!35!black,title={Before you move on}]
What does ràngzuò mean? (\textquotedblleft{}To give up one's seat (to someone)\textquotedblright{}.) Say it once more.
\end{tcolorbox}
\section[guǎngchǎngwǔ]{guǎngchǎngwǔ — square dancing}

\label{lesson:ZH-C376-guangchangwu}



\begin{quote}
Before the new one: say the Mandarin for to treat with care, then the Mandarin for to give up one's seat.
\end{quote}

\subsection*{You'll want to know: guǎngchǎngwǔ}
\emph{guǎngchǎngwǔ} — \textquotedblleft{}square dancing\textquotedblright{}.

\begin{quote}
\emph{wǒ māma měitiān wǎnshang dōu qù tiào guǎngchǎngwǔ} — My mum goes square dancing every evening.
\end{quote}

Its characters come later; for now it is a word for the ear.

Three third tones meet in this word: guǎng and chǎng are both said with a rising tone before wǔ, and the spelling keeps all three marks.

\begin{cousinweb}
\emph{guǎngchǎngwǔ} is built from \emph{guǎngchǎng} 'a public square' and \emph{wǔ} 'to dance'. All three parts are inherited from Middle Chinese, \emph{kwangX}, \emph{drjang} and \emph{mjuX} (Guangyun, in Baxter's transcription), which glosses the first 'wide' and the second 'a threshing floor'. The dancing itself is a custom of recent decades.
\end{cousinweb}

\begin{grammarlens}[title={What you've built}]
One more word for giving reasons.
\end{grammarlens}

\subsection*{Your turn}
\begin{itemize}
\raggedright
  \item \emph{Say it:} \emph{guǎngchǎngwǔ}
  \item \emph{Say it:} \emph{guǎngchǎngwǔ}, once more
  \item \emph{Say it:} \emph{wǒ māma měitiān wǎnshang dōu qù tiào guǎngchǎngwǔ} — My mum goes square dancing every evening.
  \item \emph{Recall:} say the Mandarin for to treat with care, then the Mandarin for to give up one's seat, then say \emph{guǎngchǎngwǔ} again
\end{itemize}

\begin{tcolorbox}[breakable,colback=teal!4,colframe=teal!35!black,title={Before you move on}]
What does guǎngchǎngwǔ mean? (\textquotedblleft{}Square dancing\textquotedblright{}.) Say it once more.
\end{tcolorbox}
